% Member 3 - final report, IV.E Per-Class Analysis
\subsection{Per-Class Analysis}
\label{sec:per_class}

The cross-domain average hides a great deal: two groups with the same
average can distribute their accuracy over the seven classes very
differently. Figure~\ref{fig:per_class} gives the per-class breakdown of
the main groups, and it shows where the styled groups actually win.

\paragraph{Where the gain comes from.}
The gain over the Original group is concentrated in the classes whose
identity depends on a shape rather than on a texture. \texttt{guitar} is
the clearest case: its outline --- the body, the neck, the strings --- is
what identifies it in a painting or a cartoon, and it is precisely that
outline that the stylisation threatens. \texttt{horse} and \texttt{dog}
behave similarly. The classes whose appearance is dominated by a uniform
texture move much less, because there is less shape to lose.

\paragraph{The classes that do not improve.}
\texttt{person} is the exception and deserves to be named. Its accuracy
under the styled groups is close to the Original group's and is the most
unstable across seeds. A person is defined by proportions and clothing
rather than by a compact contour, so the structure map has little to
protect: the \texttt{person} images have the lowest edge density of the
seven classes, and a modulation that is nearly uniform in effect
(Section~\ref{sec:mechanism}) cannot help them. The consistent
underperformance on \texttt{person} is also the main source of variance in
the cross-domain average, which is why the per-class view is necessary for
reading the aggregate at all.

\paragraph{The seed-level instability.}
The per-class numbers also expose a seed-level effect that the averages
hide. The styled groups' cross-domain average varies by one to two points
across the five seeds, but the variation is not spread evenly over the
classes: it is concentrated in \texttt{person} and \texttt{horse}. The
classes whose accuracy is stable across seeds are the ones whose
appearance is stable, which is consistent with the styled data being
informative but not sufficient on its own.

\paragraph{What the per-class view does not change.}
The per-class analysis does not overturn the aggregate conclusion. The
gains it shows are gains of the styled groups over the Original group, and
they are the same gains that the strength-matched control also obtains. The
styled and the matched uniform groups agree class by class as well as on
average, which is further evidence that the differences are driven by the
amount of style and not by its spatial arrangement.
